% adapted by WS from SSR's Word document and from the 5-part series at
% https://www.overleaf.com/learn/latex/How_to_Write_a_Thesis_in_LaTeX_(Part_1):_Basic_Structure
% reviewed by MAM

\documentclass[12pt,twosided]{report}

\usepackage[titletoc]{appendix} % for adding appendix to TOC
\usepackage[backend=biber]{biblatex}
\addbibresource{references.bib}
\usepackage{geometry}  % better margins and margin control
\usepackage{graphicx}  % to include figures
\usepackage{hyperref}  % internal and external links
\usepackage[utf8]{inputenc}  % support for non-ASCII characters
\usepackage{listings}  % typset code
\usepackage{outlines}  % easy nesting of lists
\usepackage{tabularx}  % more control over table column width
\usepackage{titlesec}  % customize chapter title 
\usepackage{upquote}  % prevent mishandling of single quotes in listings
\usepackage{enumitem}
\usepackage{amsmath}
\usepackage{amssymb}

% TODO: Customize the appearance of hyperref links using \hypersetup
% See https://en.wikibooks.org/wiki/LaTeX/Hyperlinks#Customization

% Custom format of chapter title.
\titleformat{\chapter}[hang]{\bf\huge}{\thechapter.}{2pc}{}

% Separate folder for images named 'images'.
\graphicspath{ {images/} }

% Separate file for references
\addbibresource{references.bib}

% Replace with your title
\title{TeachWise: An Intelligent CMS}

% Allow recalling document title
% from https://tex.stackexchange.com/a/15806/44301
\makeatletter\let\Title\@title\makeatother  

\begin{document}

\include{titlepage}  % title page.
\include{approval}  % approval page.

\chapter*{Dedication}
\textit{To our families, whose unwavering support and patience made this journey possible.}\\[1em]
\textit{To our advisor, Dr.\ Syeda Saleha Raza, whose guidance shaped the direction of this work.}\\[1em]
\textit{And to the instructors of CodeSchool, whose daily challenges inspired every line of this system.}

\chapter*{Acknowledgements}

We would like to express our deepest gratitude to our supervisor, Dr.\ Syeda Saleha Raza (Faculty of Computer Science, Habib University), for her invaluable guidance, constructive feedback, and continued support throughout both semesters of this project. Her expertise and encouragement were instrumental in shaping TeachWise into what it is today.

We are grateful to the Faculty of Computer Science at Habib University for providing us with the feedback and academic environment necessary to carry out this project.

We sincerely thank the team at CodeSchool for graciously sharing their operational data, providing access to historical scheduling records, and offering real-world context that grounded our system design and evaluation.

Finally, we thank our families and friends for their patience, motivation, and moral support throughout this demanding endeavour.

\chapter*{Abstract}

CodeSchool is an online coding education platform that connects students with instructors for live, interactive classes. As the organization has rapidly expanded to serve over 20 countries and thousands of classes, it faces significant challenges in managing instructors, class scheduling, student assignments, and payment tracking. The current manual and disconnected processes lead to inefficiencies, scheduling conflicts, and inconsistencies in instructor-student matching. 

Tracking instructor performance and maintaining centralized access to class-related data has become increasingly difficult, considering that their current teaching quality assessment relies on manual classroom observations, which are resource-intensive, subjective, and difficult to scale. The absence of objective and automated evaluation methods limits continuous feedback for instructor development and leads to inconsistencies in assessment standards, resulting in administrative overhead and limited scalability. We propose a hybrid automated system for
teaching quality assessment that combines prosodic analysis
with curriculum fidelity evaluation.

To address the issue of scheduling classes, we have also developed an intelligent scheduling algorithm that optimizes class assignments using a Genetic Algorithm. This algorithm considers multiple factors such as instructor availability, student timing preferences, and class requirements to generate optimal schedules that minimize conflicts and maximize resource utilization.

In this work, we present the design and implementation of the TeachWise platform, which integrates the intelligent scheduling algorithm and the hybrid automated teaching quality assessment system. We evaluate the effectiveness of our solutions through comprehensive testing and analysis, demonstrating significant improvements in scheduling efficiency and instructor performance evaluation compared to traditional methods.

% The following are automatically populated by LaTeX \chapter, \section and related, \figure, and \table.
\tableofcontents
\listoffigures
\listoftables

\chapter{Introduction}
\label{chap:intro}
\input{chapters/introduction}

\chapter{Literature Review}
\label{chap:lit}
\input{chapters/review}

\chapter{Software Requirement Specification (SRS)}
\label{chap:srs}
\input{chapters/srs}

\chapter{Software Design Specification (SDS)}
\label{chap:sds}
\input{chapters/sds}

\chapter{Testing and Analysis}
\label{chap:results}
% !TeX root = ../report.tex

\section{Methodology}
\label{sec:methodology}

\subsection{Evaluation}
The TeachWise evaluation framework utilizes a hybrid approach that processes classroom audio to assess both instructional delivery and semantic content coverage.

\subsubsection{Dataset and Preprocessing}
The dataset comprises 113 classroom recordings sourced from a partner online learning platform and educational YouTube content. 
\begin{itemize}
    \item \textbf{Audio Standardization:} All recordings were standardized to 16-bit PCM, 16kHz mono WAV format to ensure consistent feature extraction across samples.
    \item \textbf{Segmentation:} Audio was partitioned into 3–4 minute clips to create standardized training samples.
    \item \textbf{Transcription:} OpenAI’s Whisper ASR was utilized to generate timestamped JSON transcripts, providing the necessary input for semantic analysis.
    \item \textbf{Ground Truth:} Experienced educators annotated the clips on a 5-point Likert scale across four dimensions: energy, patience, student attention, and lesson pace.
\end{itemize}

\subsubsection{Prosodic Quality Assessment}
This component evaluates the instructor's vocal delivery using machine learning on acoustic features.
\begin{itemize}
    \item \textbf{Feature Extraction:} 18 specific acoustic features, including pitch variability, loudness, and jitter, were extracted using Librosa and OpenSMILE.
    \item \textbf{Classification Model:} Four independent Random Forest classifiers were trained, one for each quality dimension, configured with 100 decision trees and a maximum depth of 4.
\end{itemize}

\subsubsection{Curriculum Fidelity Evaluation}
The semantic component employs a Retrieval-Augmented Generation (RAG) framework to compare lesson delivery against a predefined curriculum.
\begin{itemize}
    \item \textbf{RAG Framework:} The system utilizes LLaMA3 to perform semantic comparisons between transcripts and curriculum documents extracted directly from PDF teaching guides.
    \item \textbf{Scoring Rubric:} The model assigns a score (1–5) based on a structured rubric that assesses objective coverage and the inclusion of supplementary content.
\end{itemize}

\subsection{Scheduling}

To tackle scheduling challenges, we have developed an intelligent scheduling algorithm based on a Genetic Algorithm. It considers multiple factors to generate optimal class assignments that minimize conflicts and maximize resource utilization.

\subsubsection{Problem Formulation}
  The scheduling problem was modelled as a combinatorial optimisation task. 
  Given a set of student class requests and a set of instructors, the goal  
  is to assign each request to an instructor and time slot such that hard   
  constraints are satisfied and overall schedule quality is maximised.      
                                                     
  Each request specifies a subject (e.g. Python, Mathematics, Science), a class  
  type (one-to-one, group, custom), the number of monthly sessions
  required, and a list of preferred time slots. Each instructor carries a set of subjects they can teach and a weekly availability map.

\subsubsection{Scheduling Patterns}
  Requests are first classified into one of four scheduling patterns:       
  \begin{itemize}
    \item Weekly — $N$ sessions per month at the same repeating (day, time) slot. The standard pattern for recurring classes.
    \item Consecutive — sessions on back-to-back calendar days at the same time.  
    \item Custom — the student specifies exact required (day, time) slots that must be honoured.   
    \item Flexible — N sessions distributed across any N distinct available slots, for non-standard frequencies.     
  \end{itemize}

\subsubsection{Constraint Model}

  Constraints are divided into three tiers by severity:        
  \begin{itemize}
    \item \textbf{Hard constraints} (-50,000 per violation) must be strictly enforced to ensure a feasible schedule. These include no instructor double-booking (with an exception allowing group classes of the same subject to be merged and share a slot), no assignment outside an instructor's available hours, and custom-pattern sessions must land on their required slots.
    \item \textbf{Medium constraints} (-500 per violation) should be satisfied to maintain quality, such as ensuring the assigned instructor is qualified to teach the requested subject.
    \item \textbf{Soft constraints} contribute to the overall schedule quality score and include factors like the pedagogical match score from an ML model (+100 x score, 0-1 scale), student preferred time slot (+500), Zoom licence overflow for more than X simultaneous classes (-20 per excess), and missing or extra sessions (-2,000 per session).
  \end{itemize}             

\subsubsection{Genetic Algorithm}
A genetic algorithm (GA) was used because the search space is too large for exhaustive search and the multi-objective constraint structure makes gradient-based methods unsuitable. The GA was configured with a population of 200, 100 generations, a mutation rate of 0.10, tournament size of 5, and an elite size of 5. 
The algorithm iteratively evolves the population of schedules by selecting the fittest individuals based on the defined constraints and quality metrics, applying crossover and mutation operations to generate new candidate schedules, and retaining the best solutions across generations to converge towards an optimal schedule.

\subsubsection{Benchmarking}
  Three months of historical data (November 2024 - January 2025) from our partner instituation were parsed from operational Excel files into the GA's input format. The   GA was evaluated on the January 2025 dataset (47 requests, 7 instructors) against the actual human assignments as ground truth, across five metrics:    
  \begin{itemize}
    \item \textbf{Scheduling Coverage:} The percentage of requests that were successfully scheduled without hard constraint violations.
    \item \textbf{Student Preference Satisfaction:} The percentage of scheduled sessions that were assigned to a student-preferred time slot.
    \item \textbf{Ground Truth Alignment:} The percentage of GA-generated assignments that matched the actual human-assigned instructor and time slot.
    \item \textbf{Pedagogical Match Quality:} The average ML-predicted pedagogical match score for the assigned instructor-subject pairs.
    \item \textbf{Constraint Violations:} The total number of hard and medium constraint violations in the generated schedule.
  \end{itemize}

\section{Experiments and Results}
\label{sec:results}

\subsection{Evaluation}
\subsubsection{Prosodic Assessment Performance}
The Random Forest classifiers demonstrated strong predictive performance across all four teaching quality dimensions. The results on the held-out test set are summarized in Table \ref{tab:results}.

\begin{table}[h]
\centering
\caption{Model Performance on Test Set ($N=23$)}
\label{tab:results}
\begin{tabular}{|l|c|c|c|}
\hline
\textbf{Metric} & \textbf{Test Acc. (\%)} & \textbf{F1-Score} & \textbf{CV Acc. (\%)} \\ \hline
Student Attention & 87.0 & 0.879 & 91.2 \\ \hline
Energy            & 78.3 & 0.750 & 77.8 \\ \hline
Lesson Pace       & 78.3 & 0.788 & 89.3 \\ \hline
Patience          & 73.9 & 0.763 & 89.3 \\ \hline
\end{tabular}
\end{table}

\begin{itemize}
    \item \textbf{Analysis:} Student Attention achieved the highest accuracy (87.0\%) and demonstrated robust generalization stability.
    \item \textbf{Error Patterns:} Confusion matrices indicated that misclassifications typically occurred between adjacent score levels rather than extreme deviations.
    \item \textbf{Feature Importance:} Equivalent sound level emerged as the dominant predictor across three of the four dimensions, highlighting vocal intensity as a fundamental indicator of quality.
\end{itemize}

\subsubsection{Curriculum Fidelity Results}
The LLM-based RAG component successfully identified pedagogical gaps and generated structured evaluations.
\begin{itemize}
    \item \textbf{Qualitative Consistency:} The system distinguished between complete objective achievement and partial delivery, assigning intermediate scores (2–3) when concepts were demonstrated without explicit explanation.
    \item \textbf{Actionable Feedback:} The model effectively generated objective coverage statements (e.g., "3 out of 4 objectives achieved") to support instructor development.
\end{itemize}

\subsubsection{Hybrid Score Computation}
The final teaching effectiveness score for each session is computed as a composite unweighted average:
\begin{equation}
    \mathrm{Score} = \frac{E+P+A+L+C}{5}
\end{equation}
where $E$ is Energy, $P$ is Patience, $A$ is Student Attention, $L$ is Lesson Pace, and $C$ is Curriculum Fidelity.

\subsection{Scheduling}

\subsubsection{Coverage}
The GA scheduled 100\% of requests (47/47) across all subjects, with no unscheduled students.    

\subsubsection{Constraint Satisfaction}
The GA produced zero hard-constraint violations compared to the human  schedule:
\begin{table}[h]
\centering
\caption{Violation Comparison}
\label{tab:violations}
\begin{tabular}{|l|c|c|}
\hline
\textbf{Violation Type} & \textbf{GA} & \textbf{Human} \\ \hline
Availability violations & 0 & 1 \\ \hline
Skill mismatches        & 0 & 6 \\ \hline
\end{tabular}
\end{table}

The human schedule contained six instances of an instructor assigned  outside their expertise and one slot assignment outside their available hours. The GA's repair mechanism eliminated all such violations.

\subsubsection{Student Preference Satisfaction}
94.7\% of the 38 preference-eligible requests were assigned a slot matching the student's stated preferences. The human scheduler satisfied 100\% of preferences but did so at the cost of six skill mismatches and one availability violation, suggesting that human assignments traded constraint compliance for preference satisfaction in some cases.

\subsubsection{Pedagogical Match Quality}
The GA's instructor assignments achieved an average pedagogical match  score of 0.881 versus 0.718 for the human schedule — an improvement of +0.163 ($\sim$22.7\%). This indicates that including the match score  as a soft constraint in the fitness function produces meaningfully better instructor-student pairings than manual assignment.

\subsubsection{Load Distribution}
Instructor workload was slightly more concentrated in the GA output   (standard deviation 3.37 sessions) versus the human schedule (2.42). The GA naturally favours instructors with broader availability and higher match scores. Explicit fairness balancing was not included in this version and is a candidate for future work.

\chapter{Conclusion and Future Work}
\label{chap:outro}
\section{Conclusion and Future Work}
\label{sec:conclusion}

\subsection{Conclusion}
The TeachWise evaluation framework demonstrates the efficacy of a hybrid approach that unifies prosodic analysis with grounded semantic evaluation. By achieving test accuracies between 73.9\% and 87.0\%, the system proves that machine learning models can effectively approximate complex instructional dimensions such as energy, patience, and student attention. The integration of a PDF-driven RAG framework ensures that curriculum fidelity is assessed with high specificity, addressing the critical limitations of manual observation in large-scale online platforms. TeachWise thus provides a scalable, objective, and actionable solution for managing teaching quality in distributed environments.

\subsection{Future Work}
Future research will focus on several key areas to enhance the system's capabilities:
\begin{itemize}
    \item \textbf{Linguistic Diversity:} Expanding the dataset to include a broader range of academic subjects to improve model generalization.
    \item \textbf{Annotation Rigor:} Utilizing a larger pool of expert annotators to establish formal inter-rater reliability coefficients.
    \item \textbf{Multimodal Expansion:} Incorporating visual cues, such as instructor gestures and student reactions, to provide a holistic assessment of the teaching-learning process.
\end{itemize}

\chapter{Reflection}
\label{chap:Reflection}
\section{Reflections}
\label{sec:reflections}

\subsection{Individual Reflections}

\subsubsection{Aina Shakeel}

Working on TeachWise over two semesters was one of the most formative experiences of my degree. My primary responsibilities were the scheduling algorithm and the instructor evaluation pipeline, and the journey of building both from scratch taught me far more than coursework alone could have.

One of the most important lessons I took away was how industry partnerships can change unexpectedly. At the outset, our partner institution had agreed to provide us with access to labelled curriculum and evaluation data, which was central to training and scaling the curriculum fidelity model. As the project progressed, that commitment was not fulfilled --- responses slowed, priorities shifted, and the data never materialised. This was frustrating, but it taught me a critical professional lesson: you cannot let an external dependency become a single point of failure. We responded by sourcing and annotating our own dataset from publicly available educational content, establishing our own labelling protocol, and justifying every design decision on the basis of what we actually had. Going forward, I will always build in contingency plans whenever a project depends on a third party.

A second major learning was about data bias and class imbalance. Once we had our annotated dataset and trained the Random Forest classifiers, the imbalance in label distribution --- where mid-range scores dominated --- became very visible in the results. Patience, the hardest dimension to infer from audio alone, achieved the lowest accuracy (73.9\%), and error analysis confirmed that the model struggled most at the extremes of the scale. Understanding how bias manifests in model performance, not just as an abstract concept but as something you can see in a confusion matrix, made it a concrete and deeply practical lesson.

On the technical side, building the prosodic feature extraction pipeline using Librosa and OpenSMILE, working with Whisper for ASR transcription, and designing the RAG-based curriculum fidelity component using LLaMA3 were all skills I developed during this project. The scheduling algorithm --- formulating the problem as a constrained combinatorial optimisation task and implementing a Genetic Algorithm from scratch --- was particularly rewarding. Watching the GA converge and outperform the historical human schedule in constraint compliance was a genuinely satisfying result.

Finally, working in a team with members who had very different strengths and working styles was itself a learning experience. Dividing responsibilities along lines of individual strength --- rather than trying to have everyone do everything --- was something I came to appreciate deeply. Knowing when to ask for help, and knowing when to offer it, made the collaboration more effective than I expected going in.

\subsubsection{Zain Hatim}

% TODO: Add individual reflection.

\subsubsection{Naaseh Sajid}

% TODO: Add individual reflection.

\subsubsection{Ifrah Chishti}

% TODO: Add individual reflection.

\subsection{Team Reflection}

TeachWise was built by a team of four, and over the course of two semesters we developed not only a functioning platform but also a clearer understanding of what effective collaboration actually requires.

Role distribution followed individual strengths fairly naturally. Zain took ownership of the front-end --- UI design, wireframing, and the implementation of the admin and instructor dashboards --- bringing a strong design sensibility to the visual side of the system. Aina led the backend work on the scheduling algorithm and the evaluation pipeline. Naaseh supported the backend and evaluation work alongside Aina, and played a key role in integrating the various components into a coherent system. Ifrah contributed to the scheduling module in the earlier phase of the project and subsequently shifted to supporting the backend and the finance-related CMS modules, ensuring those were fully functional and integrated.

Communication was primarily asynchronous through WhatsApp, supplemented by regular in-person and online meetings for alignment, debugging sessions, and milestone reviews. The meetings became especially important during integration phases, when the outputs of one component were inputs to another and synchronisation errors would surface unexpectedly.

One consistent challenge was timeline synchronisation. When a front-end screen depended on a backend endpoint that was not yet complete, or when the evaluation pipeline required preprocessed data that was still being collected, work would stall or need to be done in isolation with mock data. We managed this by defining clear interface contracts early --- agreeing on API response structures and data formats before the underlying logic was finished --- which reduced blocking dependencies significantly.

Another challenge was the uneven cognitive load at different stages. The data collection and annotation sprint, in particular, required concentrated effort from a subset of the team over a short window. Better upfront distribution of that kind of task is something we would approach differently if starting again.

Overall, the team dynamic was strong. Disagreements were resolved through discussion, decisions were documented, and there was a shared sense of ownership over the final product. The project would not have reached this level of completeness without genuine collaboration across all four members.

\subsection{Process Reflection}

\subsubsection{What Worked Well}

Dividing the system into well-defined modules --- scheduling, evaluation, and the CMS front-end --- and assigning clear ownership for each turned out to be one of the best structural decisions we made. It allowed parallel development across the team and meant that problems in one component rarely blocked work on another.

Adopting a constraint-tier model for the scheduling algorithm (hard, medium, and soft) was also effective. It gave the Genetic Algorithm a clear and principled signal to optimise against, and the results --- zero hard-constraint violations on the evaluation dataset versus seven in the human schedule --- validated that design choice.

Using Whisper for ASR transcription was a practical decision that paid off. The quality of the transcripts on our audio corpus was consistently high enough to serve as reliable input to the RAG curriculum fidelity pipeline without significant post-processing.

Regular documentation discipline, including maintaining an up-to-date SRS and SDS as the system evolved, made integration smoother and provided a useful reference when revisiting earlier design decisions.

\subsubsection{What Could Be Improved}

The most significant process gap was our reliance on the partner institution for evaluation data. We did not establish early enough what our fallback would be if that data did not arrive, and when it did not, we lost time that a contingency plan would have preserved.

The annotation process for the audio dataset was labour-intensive and carried out late in the project timeline, compressing the time available for model experimentation. In future projects of this kind, data collection and annotation should be treated as a first-class deliverable with its own timeline and milestones, not as a precursor to the ``real'' work.

The curriculum fidelity component, while functional, was not able to be evaluated quantitatively because we lacked ground-truth fidelity labels. Building a parallel annotation track --- where educators label curriculum coverage alongside the prosodic dimensions --- from the beginning would have made this component far more rigorously evaluated.

\subsection{Plans vs.\ Achievements}

\subsubsection{What Was Proposed}

The original proposal for TeachWise described a CMS platform with five integrated modules: instructor and student management, intelligent class scheduling, automated instructor evaluation (prosodic and curriculum fidelity), payment management, and an administrative dashboard. The evaluation component was initially scoped to include prosodic analysis across four quality dimensions, with the curriculum fidelity model to be trained on labelled data provided by the partner institution.

\subsubsection{What Was Achieved}

All five core modules were implemented and integrated into a deployment-ready web platform. The scheduling algorithm achieved its design goals, producing constraint-compliant schedules that outperformed historical human assignments across all key metrics. The prosodic evaluation pipeline was implemented in full, with four trained Random Forest classifiers achieving test accuracies between 73.9\% and 87.0\%. The RAG-based curriculum fidelity component was designed, implemented, and demonstrated to produce qualitatively consistent structured evaluations.

\subsubsection{Deviations and Reasons}

The principal deviation from the original plan was in the curriculum fidelity component. The partner institution had agreed to provide access to annotated lesson recordings paired with structured curriculum documents, which would have enabled us to train a supervised curriculum fidelity classification model and benchmark its accuracy formally. That data was never made available. As a result, the curriculum fidelity model operates in a zero-shot capacity using the RAG framework with LLaMA3, rather than as a fine-tuned supervised classifier. The component is functional and produces qualitatively sound outputs, but a quantitative accuracy benchmark does not yet exist for it. This is explicitly identified as a priority for future work.

A second change was the addition of the pedagogical match ML model as a soft constraint within the scheduling fitness function. This was not part of the original proposal but emerged during implementation as a natural enhancement once the evaluation pipeline was in place. Including match scores as a scheduling signal proved to be one of the more impactful design decisions, producing a 22.7\% improvement in average instructor--student pairing quality over the human baseline.

These deviations were not failures of planning but rather honest responses to the constraints and opportunities that emerged during real-world development. The inability to obtain partner data was a genuine external challenge; the addition of the pedagogical match constraint was a design opportunity that we recognised and acted on. Both are reflective of the adaptive approach that characterises professional software development.

\begin{appendices}

\titleformat{\chapter}[hang]{\bf\huge}{Appendix \thechapter.}{2pc}{}
  
% This appendix is optional.
\chapter{More Math}
This appendix provides supplementary mathematical background for the key techniques employed in TeachWise.

\section*{Genetic Algorithm Fitness Function}

The fitness of a candidate schedule $S$ is computed as:
\[
f(S) = \sum_{a \in S} \Bigl[ w_{\text{hard}} \cdot \delta_{\text{hard}}(a) + w_{\text{med}} \cdot \delta_{\text{med}}(a) + w_{\text{soft}} \cdot \phi_{\text{soft}}(a) \Bigr]
\]
where $\delta_{\text{hard}}(a)$ and $\delta_{\text{med}}(a)$ are penalty indicator functions for hard and medium constraint violations respectively, and $\phi_{\text{soft}}(a)$ is the soft-constraint quality score for assignment $a$. The penalty weights are set as $w_{\text{hard}} = -50{,}000$, $w_{\text{med}} = -500$, and soft-constraint bonuses include $+500$ per satisfied student time preference and $+100 \times m$ for a pedagogical match score $m \in [0,1]$.

\section*{Prosodic Feature Extraction}

Let $x(t)$ denote a raw audio waveform sampled at 16,kHz. The fundamental frequency (pitch) $F_0(t)$ is estimated frame-by-frame using the YIN algorithm. Pitch variability is captured by the standard deviation $\sigma_{F_0}$ over a 3--4 minute segment. Loudness is computed as the root-mean-square energy:
[
$\text{RMS}(t) = \sqrt{\frac{1}{N}\sum_{n=0}^{N-1} x^2(t+n)}$
]

Additional acoustic features, including jitter (cycle-to-cycle pitch variation) and shimmer (cycle-to-cycle amplitude variation), are extracted using the OpenSMILE \texttt{eGeMAPSv02} feature set. A subset of 18 selected features per segment is used as input to the Random Forest classifiers.

\section*{Random Forest Classification}

Each of the four quality-dimension classifiers (energy, patience, student attention, and lesson pace) is implemented as a Random Forest ensemble consisting of $T = 100$ decision trees. Each tree uses the Gini impurity criterion for splitting and is constrained to a maximum depth of 4 to reduce overfitting and improve generalization. Class imbalance is handled using balanced class weights. At each split, a subset of features is selected using the $\sqrt{d}$ heuristic. Final predictions are obtained through majority voting across all trees:
\[
\hat{y} = \arg\max_{c \in \{1,2,3,4,5\}}
\sum_{t=1}^{T} \mathbb{1}\big(h_t(\mathbf{x}) = c\big)
\]

where $h_t$ represents the prediction of the $t$-th tree and $c \in {1,2,3,4,5}$ is the Likert scale label.


\section*{Curriculum Fidelity Assessment}

The curriculum fidelity score is obtained using a large language model (LLaMA3) via prompt-based evaluation. The lesson transcript and corresponding curriculum content are provided directly as input to the model, which generates a structured rubric score (1--5) based on alignment with curriculum objectives and inclusion of relevant material.

This approach does not employ a retrieval-augmented generation (RAG) pipeline; no embedding-based similarity or document retrieval is performed. Instead, the evaluation is conducted using zero-shot or few-shot prompting.

\section*{Hybrid Teaching Score}

The final composite teaching effectiveness score for a session is an unweighted average of five dimensions:
\[
\text{Score} = \frac{E + P + A + L + C}{5}
\]
where $E$ is Energy, $P$ is Patience, $A$ is Student Attention, $L$ is Lesson Pace (each from the Random Forest classifiers, scaled to 1--5), and $C$ is the RAG-derived Curriculum Fidelity score (1--5).

% This appendix is required if the data set is not fully described in the main text.
\chapter{Data}
\section*{Evaluation}
Our evaluation of instruction quality was done by analysing and evaluating the data collected from various sources, mainly from YouTube and our partner instituations class recordings, linked \hyperlink{https://drive.google.com/file/d/1jeGrXsPUNfsRL8g55DkFX7OJQ83z5T-t/view?usp=sharing}{here}. These were then labelled by trained human evaluators and used to train the various models in our system. 

\section*{Scheduling}
Our scheduling data was collected from our partner institution, which provided us with anonymized scheduling data for 112 classes requests and 10 instructors over 3 months. This dataset included information on instructor availability, student enrollments, and various constraints (e.g., instructor skillset, student requested time). We used this data to benchmark our genetic algorithm-based scheduling system.

\section*{Disclaimer}
Data sourced from CodeSchool is not publicly shareable due to confidentiality agreements and child protection policies. However, we have included detailed descriptions of the data collection and processing methods in the main text to ensure transparency and reproducibility of our results.

% This appendix is required if the code is not fully described in the main text.
\chapter{Code}
\input{chapters/appendix-code}
\end{appendices}

% Print the bibliography with a ToC entry and titled, "References".

% \printbibliography[heading=bibintoc,title={References}]
\printbibliography
\end{document}

%%% Local Variables:
%%% mode: latex
%%% TeX-master: t
%%% End:
